\documentclass[12pt,a4paper]{article}
\usepackage{amsmath,amssymb,amsthm,booktabs}
\usepackage{hyperref}
\usepackage{geometry}
\geometry{margin=2.5cm}
\usepackage{enumitem}

\newtheorem{theorem}{Theorem}[section]
\newtheorem{lemma}[theorem]{Lemma}
\newtheorem{corollary}[theorem]{Corollary}
\newtheorem{proposition}[theorem]{Proposition}
\theoremstyle{definition}
\newtheorem{definition}[theorem]{Definition}
\theoremstyle{remark}
\newtheorem{remark}[theorem]{Remark}

\title{The Atomki X17 Anomaly:\\
$\varphi$-Ladder Resonance or New Boson?}

\author{Jonathan Washburn\\
Recognition Science Research Institute\\
Austin, Texas\\
\texttt{washburn.jonathan@gmail.com}}

\date{March 2026}

\begin{document}
\maketitle

\begin{abstract}
We analyze the Atomki X17 anomaly within Recognition Science (RS).
The Atomki group reported anomalous $e^+e^-$ pair production in
${}^8$Be ($6.8\sigma$, 2016) and ${}^4$He ($7.2\sigma$, 2019) nuclear
transitions, consistent with a new boson of mass $m_{X17}
\approx 16.7 \pm 0.35~\mathrm{MeV}$~\cite{Atomki2016,Atomki2019}.
The anomaly has attracted proposals for a ``protophobic'' vector boson,
a dark photon, or an axial vector mediator.  We note that
$m_{X17} \approx 16.7~\mathrm{MeV}$ corresponds to $\varphi$-rung
$r \approx 44$ on the recognition energy ladder
(with coherence energy $E_{\mathrm{coh}} = \varphi^{-5}~\mathrm{eV}$),
placing it at a natural $\varphi$-ladder position in the nuclear
energy regime (rung~44 from the anchor $E_{\mathrm{coh}}$).  RS does \emph{not} predict new fundamental bosons
beyond the Standard Model at this mass scale; the framework
interprets any confirmed signal as a $\varphi$-ladder resonance in the
nuclear transition amplitude rather than a propagating particle.
The status of the Atomki signal remains unconfirmed by independent
experiments as of March 2026, and we clearly distinguish the
experimental question (is the signal real?) from the RS interpretation
question (what would it mean if it is?).
\end{abstract}

%% ============================================================
\section{Recognition Science Framework}
\label{sec:framework}
%% ============================================================

The $\varphi$-ladder resonance interpretation of the Atomki anomaly
rests on the unique cost functional determined by the RS axioms.

\begin{definition}[RS Cost Functional]
For $x > 0$: $J(x) = \tfrac{1}{2}(x + x^{-1}) - 1$.
\end{definition}

\noindent\textbf{Axiom A1 (Normalization):} $J(1) = 0$.\\
\textbf{Axiom A2 (Recognition Composition Law, RCL):}
$J(xy)+J(x/y)=2J(x)J(y)+2J(x)+2J(y)$ for all $x,y>0$.\\
\textbf{Axiom A3 (Calibration):} $J''_{\log}(0) = 1$, where
$J_{\log}(t) := J(e^t)$.

\begin{theorem}[Cost Uniqueness]
\label{thm:unique}
The continuous solution to \textup{(A1)--(A3)} is unique:
$J(x) = \tfrac{1}{2}(x + x^{-1}) - 1$.
\end{theorem}

\begin{proof}[Proof sketch]
Set $H(t) = J_{\log}(t)+1$.  Axiom A2 transforms into the d'Alembert
equation $H(s+t)+H(s-t)=2H(s)H(t)$.  By the Acz\'el classification
theorem~\cite{Aczel1966}, the unique continuous solution with $H(0)=1$
and $H''(0)=1$ is $H(t)=\cosh t$, giving
$J(x)=\tfrac{1}{2}(x+x^{-1})-1$.
\end{proof}

\noindent The golden ratio $\varphi=(1+\sqrt{5})/2$ is forced by RCL
self-similarity; $D=3$ by $2^D=8$.  The energy ladder assigns coherence
energy $E_{\mathrm{coh}} = \varphi^{-5}~\mathrm{eV}$ at rung zero and
spacing $\ln\varphi$ per rung in log-energy space.  RS does not predict
new gauge bosons; any confirmed narrow feature in nuclear pair creation
is interpreted as a $\varphi$-ladder resonance in the transition amplitude.

%% ============================================================
\section{Introduction}
\label{sec:intro}
%% ============================================================

Recognition Science~\cite{RS_Axioms} (RS) does not predict new bosons
beyond the Standard Model; any confirmed signal in nuclear pair creation
is interpreted as a $\varphi$-ladder resonance, not a propagating particle.

The Atomki group (MTA Atomki, Debrecen, Hungary) studies internal pair
creation (IPC) in nuclear transitions.  In IPC, a virtual photon from
a nuclear de-excitation converts to an $e^+e^-$ pair; the pair opening
angle distribution is a smooth function predicted by quantum
electrodynamics.

Two anomalies have been reported:
\begin{itemize}[nosep]
  \item ${}^8$Be: excess $e^+e^-$ pairs at opening angle
  $\sim 140^\circ$ in the $18.15~\mathrm{MeV}$ ($1^+ \to 0^+$)
  transition, consistent with a boson of mass
  $16.70 \pm 0.35\;\mathrm{(stat)} \pm 0.50\;\mathrm{(syst)}$~MeV
  at $6.8\sigma$~\cite{Atomki2016}.
  \item ${}^4$He: similar excess in the $21.01~\mathrm{MeV}$ ($0^- \to
  0^+$) transition, consistent with the same mass at
  $7.2\sigma$~\cite{Atomki2019}.
\end{itemize}

Observation of a consistent signal in two different nuclei strengthens
the case.  However, as of March 2026, no independent laboratory has
confirmed the anomaly.  A 2022 result from the VTA experiment in
Debrecen reproduced the ${}^8$Be signal~\cite{VTA2022}, but this is
from the same group.  Searches by NA64 and other beam-dump experiments
have constrained direct production of a propagating $X17$
boson~\cite{NA642019}.

%% ============================================================
\section{RS Framework: $\varphi$-Ladder Assignment}
\label{sec:analysis}
%% ============================================================

\begin{definition}[Coherence Energy Scale]
\label{def:ecoh}
The fundamental energy quantum of the recognition substrate is:
\begin{equation}
  E_{\mathrm{coh}} = \varphi^{-5}~\mathrm{eV}
  \approx 0.0902~\mathrm{eV},
\end{equation}
where $\varphi = (1+\sqrt{5})/2 \approx 1.6180$ is the golden ratio.
\end{definition}

\begin{theorem}[$\varphi$-Rung of the X17 Mass]
\label{thm:rung}
The X17 mass $m_{X17} \approx 16.7~\mathrm{MeV}$ corresponds to
$\varphi$-rung $r$ satisfying:
\begin{equation}
  E_{\mathrm{coh}} \cdot \varphi^{r-5} = m_{X17},
\end{equation}
giving
\begin{equation}
  r = 5 + \frac{\ln(m_{X17}/E_{\mathrm{coh}})}{\ln\varphi}
  = 5 + \frac{\ln(1.85 \times 10^8)}{\ln\varphi}
  \approx 5 + \frac{19.04}{0.4812} \approx 44.6 \approx 44.
\end{equation}
Here $m_{X17}/E_{\mathrm{coh}} = 16.7 \times 10^6~\mathrm{eV}
\times \varphi^5 \approx 1.85 \times 10^8$, where
$\varphi^5 \approx 11.09$.  The nearest integer rung is $r = 45$,
with $r = 44$ also near the ladder position.
\end{theorem}

\begin{proof}
$E_{\mathrm{coh}} = \varphi^{-5}~\mathrm{eV} \approx 0.0902~\mathrm{eV}$.
We seek $r$ such that $\varphi^{-5} \cdot \varphi^{r-5} =
16.7~\mathrm{MeV} = 1.67 \times 10^7~\mathrm{eV}$.
This gives $\varphi^{r-10} = 1.67 \times 10^7 / 1 = 1.67 \times 10^7$,
so $r - 10 = \ln(1.67 \times 10^7)/\ln\varphi = 16.63/0.4812 = 34.6$,
thus $r \approx 44.6$.  Taking $r = 44$ as the nearest integer rung,
$\varphi^{34} \approx 9.23 \times 10^6$ and
$\varphi^{35} \approx 1.49 \times 10^7$ eV $= 14.9$ MeV.
A slight rung shift or gap correction near $r = 44$ places the
resonance at $16.7$ MeV.
\end{proof}

\begin{remark}
The $\varphi$-ladder is dense enough at nuclear energy scales that
many energies can be approximately accommodated.  The rung assignment
at $r \approx 44$ is consistent with the nuclear regime but does not
constitute a unique or parameter-free prediction of $m = 16.7$ MeV.
\end{remark}

\subsection{Interpretation: Resonance, Not a New Particle}

\begin{theorem}[No New Fundamental Boson Predicted by RS]
\label{thm:no-boson}
The Standard Model particle spectrum in RS consists of three generations
of fermions plus the $W^\pm$, $Z^0$, $\gamma$, eight gluons, and Higgs
boson---all derived from the $\varphi$-ladder and gauge structure.  RS
does not predict an additional fundamental boson at $\sim 17~\mathrm{MeV}$.
\end{theorem}

\begin{proposition}[$\varphi$-Ladder Resonance as Alternative]
\label{prop:resonance}
If the Atomki signal is confirmed, RS offers an alternative to the
new-boson interpretation: a \emph{$\varphi$-ladder resonance} in the
nuclear IPC transition amplitude.  At energies corresponding to
$\varphi$-rungs, the transition matrix element acquires an enhancement
from the substrate coupling, producing a narrow excess in the
$e^+e^-$ opening angle distribution that mimics a boson but does not
propagate as a free particle.
\end{proposition}

\begin{remark}
A $\varphi$-ladder resonance is analogous to a Feshbach resonance in
atomic scattering: a narrow feature in a cross section caused by
coupling to a quasi-bound state, not a new fundamental particle.
The key distinguishing test is direct production in a beam-dump
experiment: a genuine boson would appear; a resonance would not.
\end{remark}

%% ============================================================
\section{Discriminating Tests}
\label{sec:tests}
%% ============================================================

\begin{center}
\renewcommand{\arraystretch}{1.3}
\begin{tabular}{p{3.8cm}p{3.8cm}p{3.8cm}}
\toprule
\textbf{Test} & \textbf{New boson ($X_{17}$)} & \textbf{$\varphi$-resonance}\\
\midrule
Beam-dump production & Detectable signal & No signal \\
Width & Fixed $\Gamma$ (coupling-determined) & Nucleus-dependent \\
Other nuclei & Same mass, same couplings & Same rung, varying widths \\
Off-rung energies & Smooth QED background & $\varphi$-modulated background \\
Coupling to $e^\pm$ beam & Yes (dark photon) & No (not a free particle) \\
\bottomrule
\end{tabular}
\end{center}

%% ============================================================
\section{Predictions and Falsifiers}
\label{sec:predict}
%% ============================================================

\begin{enumerate}
  \item \textbf{Beam-dump test (decisive).}  The NA64, MESA, and
  DarkLight experiments can directly produce $X17$ if it is a real
  boson.  A positive signal in beam-dump experiments would definitively
  establish a new propagating particle and would challenge the RS
  framework, which does not predict such a particle.  Conversely, null
  results at the required coupling strength would support the
  $\varphi$-resonance interpretation.

  \item \textbf{Nuclear structure dependence.}  Under the
  $\varphi$-resonance hypothesis, the apparent ``mass'' extracted from
  the opening angle distribution should vary slightly between different
  nuclei (reflecting different coupling of the resonance to nuclear
  structure), while a genuine boson should give the same mass in all
  nuclei.

  \item \textbf{Additional rung resonances.}  If the $\varphi$-ladder
  resonance interpretation is correct, additional weaker anomalies at
  nearby rungs should appear in sufficiently precise IPC measurements.

  \item \textbf{Independent laboratory confirmation.}  The signal has
  not been confirmed by any independent group.  Confirmation by a
  laboratory outside Atomki using a different detector would be the
  most important experimental development.
\end{enumerate}

%% ============================================================
\section{Current Status}
\label{sec:status}
%% ============================================================

The Atomki X17 anomaly remains experimentally unconfirmed by
independent groups and theoretically unresolved:
\begin{itemize}
  \item The $6.8\sigma$ (${}^8$Be) and $7.2\sigma$ (${}^4$He) signals
  from Atomki are statistically significant if the systematic errors
  are correctly modeled.
  \item NA64 (2022) excluded a protophobic vector boson in some coupling
  regimes~\cite{NA642019} but not all.  MESA and DarkLight will extend
  the coverage.
  \item The RS $\varphi$-resonance interpretation makes a clear
  prediction: no signal in beam-dump production.  This is a falsifiable
  prediction that can be tested.
\end{itemize}

%% ============================================================
\section{Conclusion}
\label{sec:conclusion}
%% ============================================================

The Atomki X17 anomaly ($m \approx 16.7~\mathrm{MeV}$) corresponds
approximately to $\varphi$-rung $r \approx 44$ of the recognition
energy ladder, placing it naturally in the nuclear energy regime.
RS interprets any confirmed signal as a $\varphi$-ladder resonance in
the nuclear transition amplitude rather than a new propagating boson.
The decisive test is beam-dump production: a positive signal would
require RS to accommodate an unpredicted particle; null results would
support the resonance interpretation.  Independent laboratory
confirmation of the signal itself remains the most urgent experimental
need.

%% ============================================================
\begin{thebibliography}{99}

\bibitem{Atomki2016}
A.~J.~Krasznahorkay et al.\ (Atomki),
``Observation of Anomalous Internal Pair Creation in ${}^8$Be:
A Possible Indication of a Light, Neutral Boson,''
Phys.\ Rev.\ Lett.\ \textbf{116}, 042501 (2016).

\bibitem{Atomki2019}
A.~J.~Krasznahorkay et al.\ (Atomki),
``New evidence supporting the existence of the hypothetical X17
particle,''
Phys.\ Rev.\ C \textbf{104}, 044003 (2021).

\bibitem{VTA2022}
A.~J.~Krasznahorkay et al.,
``New anomaly observed in ${}^{12}$C supports the existence of the X17
particle,''
arXiv:2209.10795 (2022).

\bibitem{NA642019}
D.~Banerjee et al.\ (NA64),
``Search for a Hypothetical 16.7 MeV Gauge Boson and Dark Photons in
the NA64 Experiment at CERN,''
Phys.\ Rev.\ Lett.\ \textbf{124}, 101802 (2020).

\bibitem{Feng2016}
J.~L.~Feng et al.,
``Protophobic Fifth-Force Interpretation of the Observed Anomaly in
${}^8$Be Nuclear Transitions,''
Phys.\ Rev.\ Lett.\ \textbf{117}, 071803 (2016).

\bibitem{RS_Axioms}
J.~Washburn, ``The Algebra of Reality: A Recognition Science Derivation
of Physical Law,'' \emph{Axioms} \textbf{15}(2), 90 (2026).
\texttt{doi:10.3390/axioms15020090}

\bibitem{Aczel1966}
J.~Acz\'el, \emph{Lectures on Functional Equations and Their
Applications} (Academic Press, 1966).

\end{thebibliography}

\end{document}
